%%%%%%%%%%%%%%%%%%%%%%%%%%%%%%%%%%%%%%%%%
% Beamer Presentation
% LaTeX Template
% Version 1.0 (10/11/12)
%
% This template has been downloaded from:
% http://www.LaTeXTemplates.com
%
% License:
% CC BY-NC-SA 3.0 (http://creativecommons.org/licenses/by-nc-sa/3.0/)
%
%%%%%%%%%%%%%%%%%%%%%%%%%%%%%%%%%%%%%%%%%

%----------------------------------------------------------------------------------------
%	PACKAGES AND THEMES
%----------------------------------------------------------------------------------------

\documentclass{beamer}

\mode<presentation> {

% The Beamer class comes with a number of default slide themes
% which change the colors and layouts of slides. Below this is a list
% of all the themes, uncomment each in turn to see what they look like.

%\usetheme{default}
\usetheme{AnnArbor}
%\usetheme{Antibes}
%\usetheme{Bergen}
%\usetheme{Berkeley}
%\usetheme{Berlin}
%\usetheme{Boadilla}
%\usetheme{CambridgeUS}
%\usetheme{Copenhagen}
%\usetheme{Darmstadt}
%\usetheme{Dresden}
%\usetheme{Frankfurt}
%\usetheme{Goettingen}
%\usetheme{Hannover}
%\usetheme{Ilmenau}
%\usetheme{JuanLesPins}
%\usetheme{Luebeck}
%\usetheme{Madrid}
%\usetheme{Malmoe}
%\usetheme{Marburg}
%\usetheme{Montpellier}
%\usetheme{PaloAlto}
%\usetheme{Pittsburgh}
%\usetheme{Rochester}
%\usetheme{Singapore}
%\usetheme{Szeged}
%\usetheme{Warsaw}

% As well as themes, the Beamer class has a number of color themes
% for any slide theme. Uncomment each of these in turn to see how it
% changes the colors of your current slide theme.

%\usecolortheme{albatross}
%\usecolortheme{beaver}
%\usecolortheme{beetle}
%\usecolortheme{crane}
%\usecolortheme{dolphin}
%\usecolortheme{dove}
%\usecolortheme{fly}
%\usecolortheme{lily}
%\usecolortheme{orchid}
%\usecolortheme{rose}
%\usecolortheme{seagull}
%\usecolortheme{seahorse}
%\usecolortheme{whale}
%\usecolortheme{wolverine}

%\setbeamertemplate{footline} % To remove the footer line in all slides uncomment this line
%\setbeamertemplate{footline}[page number] % To replace the footer line in all slides with a simple slide count uncomment this line

%\setbeamertemplate{navigation symbols}{} % To remove the navigation symbols from the bottom of all slides uncomment this line
}

\usepackage{graphicx} % Allows including images
\usepackage{booktabs} % Allows the use of \toprule, \midrule and \bottomrule in tables

%\usepackage{ctex}

%----------------------------------------------------------------------------------------
%	TITLE PAGE
%----------------------------------------------------------------------------------------

\title[Numerical Analysis]{Team 5 Presentation} % The short title appears at the bottom of every slide, the full title is only on the title page

\author{Team 5} % Your name
\institute[QUT-Math] % Your institution as it will appear on the bottom of every slide, may be shorthand to save space
{
Qingdao University of Technology \\ % Your institution for the title page
\medskip
\textit{} % Your email address
}
\date{\today} % Date, can be changed to a custom date
%------------------------------------------------------------------------
\begin{document}
\begin{frame}
	\titlepage % Print the title page as the first slide
\end{frame}
%------------------------------------------------------------------------
\begin{frame}
\frametitle{放假啊点上课呢热搜刚回家啊虽然离开嘎的时间管理的撒反馈；加管理的撒开}
由胡克定律可得经典谐振子模型
\begin{equation}
F=-kx=m\frac{d^2x}{dt^2}
\end{equation}
其解为
\begin{equation}
x(t)=Asin(\omega t)+Bcos(\omega t)
\end{equation}
其中
\begin{equation}
\omega \equiv \sqrt{\frac{k}{m}}
\end{equation}
为谐振子的角频率，可得势能为
\begin{equation}
V(x)=\frac{1}{2}m\omega^2 x^2
\end{equation}
同时上式也为量子力学问题要求解的势能，即求解的薛定谔方程为
\begin{equation}
-\frac{\hbar^2}{2m}\frac{d^2\psi}{dx^2}+\frac{1}{2}m\omega^2x^2\psi=E\psi
\end{equation}
\end{frame}
%------------------------------------------------------------------------
\begin{frame}
	\frametitle{目录} % Table of contents slide, comment this block out to remove it
	\tableofcontents % Throughout your presentation, if you choose to use \section{} and \subsection{} commands, these will automatically be printed on this slide as an overview of your presentation
\end{frame}
\section{代数法}
%--------------------------------------------------------------------
\begin{frame}
\frametitle{代数法}
将(5)式写成下面的形式
\begin{equation}
\frac{1}{2m}\left[ p^2+(m\omega x)^2\right]\psi=E\psi 
\end{equation}
如何求解上面的方程，基本的思想就是分解哈密顿算符
\begin{equation}
H=\frac{1}{2m}\left[ p^2+(m\omega x)^2\right] 
\end{equation}
引入以下变量
\begin{equation}
a_{\pm}\equiv\frac{1}{\sqrt{2\hbar m\omega}}(\mp ip+m\omega x)
\end{equation}

\end{frame}
%--------------------------------------------------
\begin{frame}
计算$a_{-}a_{+}$的乘积
\begin{equation}
\begin{aligned}
a_{-}a_{+}
&=\frac{1}{2\hbar m\omega}(ip+m\omega x)(-ip+m\omega x)\\
&=\frac{1}{2\hbar m\omega}\left[ p^2+(m\omega x)^2-im\omega(xp-px)\right] 
\end{aligned}
\end{equation}
其中有一个额外的项，涉及到$(xp-px)$,我们称其为x与p的对易式。一般情况，算符A和B的对易式是
\begin{equation}
\left[ A,B\right]=AB-BA 
\end{equation}
这样，(9)式可表示为
\begin{equation}
a_{-}a_{+}=\frac{1}{2\hbar m\omega}\left[ p^2+(m\omega x)^2\right] -\frac{i}{2\hbar}\left[ x,p\right]  
\end{equation}

\end{frame}
%---------------------------------------------------
\begin{frame}
求解x和p的对易式，引入测试函数f(x)
\begin{equation}
\left[ x,p\right] f(x)=\left[ x\frac{\hbar}{i}\frac{d}{dx}(f)-\frac{\hbar}{i}\frac{d}{dx}(xf)\right]=\frac{\hbar}{i}(x\frac{df}{dx}-x\frac{df}{dx}-f)=i\hbar f(x) 
\end{equation}
最终可得到
\begin{equation}
\left[ x,p\right] =i\hbar
\end{equation}
上式称为正则对易关系\\
利用该关系，(11)式可写为
\begin{equation}
a_{-}a_{+}=\frac{1}{\hbar\omega}H+\frac{1}{2}
\end{equation}
若$a_{+}$在左边，则有
\begin{equation}
a_{+}a_{-}=\frac{1}{\hbar\omega}H-\frac{1}{2}
\end{equation}
由以上两式可得
\begin{equation}
\left[ a_{-},a_{+}\right] =1
\end{equation}

\end{frame}
%--------------------------------------------------
\begin{frame}
根据(14)(15)式，我们可以将哈密顿量写为
\begin{equation}
H=\hbar\omega(a_{-}a_{+}-\frac{1}{2})
\end{equation}
和
\begin{equation}
H=\hbar\omega(a_{+}a_{-}+\frac{1}{2})
\end{equation}
利用$a_{\pm}$,可以将谐振子的薛定谔方程写成如下的形式
\begin{equation}
\hbar\omega(a_{\pm}a_{\mp}\pm\frac{1}{2})\psi=E\psi
\end{equation}
\end{frame}
\begin{frame}

若$\psi$能够满足能量为E的薛定谔方程$H\psi=E\psi$,则$a_{+}\psi$就会满足能量为$(E+\hbar\omega)$的薛定谔方程$H(a_{+}\psi)=(E+\hbar\omega)(a_{+}\psi)$，证明如下
\begin{equation}
\begin{aligned}
H(a_{+}\psi)
&=\hbar\omega(a_{+}a_{-}+\frac{1}{2})(a_{+}\psi)\\
&=\hbar\omega(a_{+}a_{-}a_{+}+\frac{1}{2}a_{+})\psi\\
&=\hbar\omega a_{+}(a_{-}a_{+}+\frac{1}{2})\psi\\
&=a_{+}\left[ \hbar\omega(a_{+}a_{-}+1+\frac{1}{2})\psi\right] \\
&=a_{+}(H+\hbar\omega)\psi=a_{+}(E+\hbar)\psi=(E+\hbar\omega)(a_{+}\psi)
\end{aligned}
\end{equation}
同理可证，$a_{-}\psi$是能量为$(E-\hbar\omega)$的解。
\end{frame}
%--------------------------------------------------------
\begin{frame}
在上述方法中，若我们得到一个解，我们就可以通过升降能量得到其它的解，我们把$a_{\pm}$升降阶算符。\\
如果我们反复利用升降阶算符，就会得到低于零的能量状态，这种情况并不存在，所以存在一个最低的波函数使得
\begin{equation}
a_{-}\psi_0=0
\end{equation}
利用上式我们可以确定$\psi_0(x)$
\begin{equation}
\psi_0(x)=(\frac{m\omega}{\pi\hbar})^{\frac{1}{4}}e^{-\frac{m\omega}{2\hbar}x^2}
\end{equation}
将其带入(19)式，$\hbar\omega(a_{+}a_{-}+\frac{1}{2})\psi_0=E_0\psi_0$,可以得到基态能量
\begin{equation}
E_0=\frac{1}{2}\hbar\omega
\end{equation}
\end{frame}
%------------------------------------------------------------
\begin{frame}
这样我们就可以通过将升降阶算符作用于$\psi_0$,原则上能够得到谐振子的所有定态
\begin{equation*}
\psi_n(x)=A_n(a_{+})^n\psi_0(x)
\end{equation*}
\begin{equation}
E_n=(n+\frac{1}{2})\hbar\omega
\end{equation}
接下来求解归一化常数$A_n$
\end{frame}
%---------------------------------------------------------------
\begin{frame}
根据上式，我们可以得知$a_{\pm}$正比于$\psi_{n\pm 1}$,即
\begin{equation}
a_{+}\psi_n=c_n\psi_{n+1}\quad,\quad a_{-}\psi_n=d_n\psi_{n-1}
\end{equation}
注意到，对于任何可积函数$f(x)$和$g(x)$
\begin{equation}
\int_{-\infty}^{+\infty}f^*(a_{\pm}g)dx=\int_{-\infty}^{\infty}(a_{mp}f)^*gdx
\end{equation}
即$a_{\mp}$是$a_{\pm}$的厄密共轭，证明如下
\begin{equation}
\begin{aligned}
\int_{-\infty}^{+\infty}f^*(a_{\pm}g)dx
&=\frac{1}{\sqrt{2\hbar m\omega}}\int_{-\infty}^{+\infty}f^*(\mp\hbar\frac{d}{dx}+m\omega x)gdx\\
&=\frac{1}{\sqrt{2\hbar m\omega}}\int_{-\infty}^{+\infty}\left[ (\pm\hbar\frac{d}{dx}+m\omega x)f\right]^*gdx\\ 
&=\int_{-\infty}^{+\infty}(a_{\mp}f)^*fdx
\end{aligned}
\end{equation}
\end{frame}
%-----------------------------------------------------------------
\begin{frame}
根据上式，特别有
\begin{equation}
\int_{-\infty}^{+\infty}(a_{\pm}\psi_n)^*(a_{\pm}\psi_n)dx=\int_{-\infty}^{+\infty}(a_{\mp}a_{\pm}\psi_n)^*\psi_ndx
\end{equation}
由(19)与(24)式可以得到
\begin{equation}
a_{+}a_{-}\psi_n=n\psi_n\quad,\quad a_{-}a_{+}\psi_n=(n+1)\psi_n
\end{equation}
所以
\begin{equation*}
\int_{-\infty}^{+\infty}(a_{+}\psi_n)^*(a_{+}\psi_n)dx=\left| c_n\right|^2\int_{-\infty}^{+\infty}\left| \psi_{n+1}\right|^2dx=(n+1)\int_{-\infty}^{+\infty}\left| \psi_n\right|^2 dx  
\end{equation*}
\begin{equation}
\int_{-\infty}^{+\infty}(a_{-}\psi_n)^*(a_{-}\psi_n)dx=\left| d_n\right|^2\int_{-\infty}^{+\infty}\left| \psi_{n-1}\right|^2dx=n\int_{-\infty}^{+\infty}\left| \psi_n\right|^2  dx 
\end{equation}
\end{frame}
\begin{frame}
根据上式可以得到$c_n$与$d_n$,即
\begin{equation}
a_{+}\psi_n=\sqrt{n+1}\psi_{n+1}\quad a_{-}\psi_n=\sqrt{n}\psi_{n-1}
\end{equation}
这样就有
$$\psi_1=a_{+}\psi_0$$
$$\psi_2=\frac{1}{\sqrt{2}}a_{+}\psi_1=\frac{1}{\sqrt{2}}(a_{+})^2\psi_0$$
$$\psi_3=\frac{1}{\sqrt{3}}a_{+}\psi_2=\frac{1}{\sqrt{3\times 2}}(a_{+})^3\psi_0$$
$$\psi_4=\frac{1}{\sqrt{4}}a_{+}\psi_3=\frac{1}{\sqrt{4\times3\times 2}}(a_{+})^4\psi_0$$
依此类推
\begin{equation}
\psi_{n}=\frac{1}{\sqrt{n!}}(a_{+})^n\psi_0
\end{equation}
\end{frame}
%----------------------------------------------------------------
\section{解析法}
\begin{frame}
\frametitle{解析法}
回到谐振子的薛定谔方程
\begin{equation*}
-\frac{\hbar^2}{2m}\frac{d^2\psi}{dx^2}+\frac{1}{2}m\omega^2x^2\psi=E\psi
\end{equation*}
引入无量纲变量
\begin{equation}
\xi\equiv\sqrt{\frac{m\omega}{\hbar}}x
\end{equation}
此时，薛定谔方程可以写为
\begin{equation}
\frac{d^2\psi}{d\xi^2}=(\xi^2-K)\psi
\end{equation}
其中
\begin{equation}
K\equiv\frac{2E}{\hbar\omega}
\end{equation}
\end{frame}
%------------------------------------------------
\begin{frame}
若$\xi$很大，则可以略去含K的项，即在这样的区域里
\begin{equation}
\frac{d^2\psi}{d\xi^2}\approx\xi^2\psi
\end{equation}
其近似解
\begin{equation}
\psi(\xi)\approx Ae^{-\frac{\xi^2}{2}}+Be^{\frac{\xi^2}{2}}
\end{equation}
含B的项不能归一化，所以，物理上可接受的解具有的渐进形式为
\begin{equation}
\psi(\xi)=h(\xi)e^{-\frac{\xi^2}{2}}
\end{equation}
\end{frame}
%-------------------------------------------------
\begin{frame}
对上式求导
$$\frac{d\psi}{d\xi}=(\frac{dh}{d\xi}-\xi h)e^{-\frac{\xi^2}{2}}$$
$$\frac{d^2\psi}{d\xi^2}=(\frac{d^2h}{d\xi^2}-2\xi\frac{dh}{d\xi}+(\xi^2-1)h)e^{-\frac{\xi^2}{2}}$$
薛定谔方程变为
\begin{equation}
\frac{d^2h}{d\xi^2}-2\xi\frac{dh}{d\xi}+(K-1)h=0
\end{equation}
现在来寻求上式的$\xi$幂级数解
\end{frame}
%-----------------------------------------------------
\begin{frame}
\begin{equation}
h(\xi)=a_0+a_1\xi+a_2\xi^2+\cdots=\sum_{j=0}^{\infty}a_j\xi^j
\end{equation}
对该级数逐项求导
$$\frac{dh}{d\xi}=a_1+2a_2\xi+3a_3\xi^2+\cdots+\cdots=\sum_{j=0}^{\infty}ja_J\xi^{j-1}$$
$$\frac{d^2h}{d\xi^2}=2a_2+2\times3a_3\xi+3\times4a_4\xi^2+\cdots=\sum_{j=0}^{\infty}(j+1)(j+2)a_{j+2}\xi^j$$
将这些带入(39)式，可得
\begin{equation}
\sum_{j=0}^{\infty}\left[ (j+1)(j+2)a_{j+2}-2ja_{j}+(K-1)a_j \right]\xi^j=0
\end{equation}
\end{frame}
%-------------------------------------------------------
\begin{frame}
要使上式成立，由幂级数展开的唯一性，每一个$\xi$幂次前的系数必须为零
$$(j+1)(j+2)a_{j+2}-2ja_j+(K-1)a_j=0$$
因此
\begin{equation}
a_{j+2}=\frac{(2j+1-K)}{(j+1)(j+2)}a_j
\end{equation}
这个递推公式完全等价于薛定谔方程，从$a_0$开始，可以给出所有的偶系数
$$a_2=\frac{(1-K)}{2}a_0,a_4=\frac{5-K}{12}a_2=\frac{(5-K)(1-K)}{24}a_1,\cdots$$
从$a_1$开始，可以给出所有的奇系数
$$a_3=\frac{(3-K)}{6}a_1,a_5=\frac{7-K}{20}a_3=\frac{(7-K)(3-K)}{120}a_1,\cdots$$
\end{frame}
%-------------------------------------------------------------
\begin{frame}
完整的解可以写为
\begin{equation}
h(\xi)=h_{even}(\xi)+h_{odd}(\xi)
\end{equation}
其中
$$h_{even}(\xi)\equiv a_0+a_2\xi^2+a_4\xi^4+\cdots$$
$$h_{odd}(\xi)\equiv a_1\xi+a_3\xi^3+a_5\xi^5+\cdots$$
对于非常大的j，递推公式变为
\begin{equation}
a_{j+2}\approx\frac{2}{j}a_j
\end{equation}
其近似解为
\begin{equation}
a_j\approx\frac{C}{(j/2)!}
\end{equation}
\end{frame}
%-------------------------------------------------------------------
\begin{frame}
在这种情况下将导致
\begin{equation}
h(\xi)\approx C \sum\frac{1}{\frac{j}{2}!}\xi^j\approx C\sum\frac{1}{j!}\xi^{2j}\approx e^{\xi^2}
\end{equation}
此时，根据(38)式，$\psi$的行为就是$e^{\frac{\xi^2}{2}}$,无法归一化，我们令最大的$j=n$,要得到可接受的解，要求
\begin{equation}
K=2n+1
\end{equation}
并且当$n$为偶数时，$a_1=0$,当$n$为奇数时，$a_0=0$\\
根据(35)式，我们可以知道能量必须是
$$E_n=(n+\frac{1}{2})\hbar\omega\quad n=1,2,\cdots$$
\end{frame}
%----------------------------------------------------------------------
\begin{frame}
对允许的K值,递推公式为
\begin{equation}
a_{j+2}=\frac{-2(n-j)}{(j+1)(j+2)}a_j
\end{equation}
当n=0时
$$h_0(\xi)=a_0$$
$$\psi_0(\xi)=a_0e^{-\frac{\xi^2}{2}}$$
当n=1时
$$h_1(\xi)=a_1\xi$$
$$\psi_1(\xi)=a_1\xi e^{\frac{-\xi^2}{2}}$$
当n=2时
$$h_2(\xi)=a_0(1-2\xi^2)$$
$$\psi_2(\xi)=a_0(1-2\xi^2)e^{\frac{-\xi^2}{2}}$$
\end{frame}
%--------------------------------------------------------------------
\begin{frame}
引入厄密多项式$H_n(\xi)$,下列为厄密多项式的前几项
\begin{equation}
\begin{aligned}
&H_0=1\\
&H_1=2\xi\\
&H_2=4\xi^2-2\\
&H_3=8\xi^3-12\xi\\
&H_4=16\xi^4-48\xi^2+12\\
&H_5=32\xi^5-160\xi^3+120\xi
\end{aligned}
\end{equation}

\end{frame}
\begin{frame}
厄密多项式的罗德里格公式
\begin{equation}
H_n(\xi)=(-1)^ne^{\xi^{2}}(\frac{d}{d\xi})^ne^{-\xi^{2}}
\end{equation}
厄密多项式的递推关系
\begin{equation}
H_{n+1}(\xi)=2\xi H_n(\xi)-2nH_{n-1}(\xi)
\end{equation}
此时,归一化的谐振子定态是
\begin{equation}
\psi_{n}(x)=(\frac{m\omega}{\pi\hbar})^{\frac{1}{4}}\frac{1}{\sqrt{2^nn!}}H_n(\xi)e^{-\frac{\xi^2}{2}}
\end{equation}
\end{frame}

\end{document}
